\chapter{Overzicht}
\section{Samenvatting van het project}
\subsection{Doel, scope en objectieven}
Het doel van dit project is het maken van een webapplicatie die het mogelijk maakt om lessenroosters binnen de universiteit te schedulen. Deze lessenroosters moeten vervolgens door de studenten geraadpleegd kunnen worden. Er is verder specifieke support voor mobiele platformen nodig.
\\
\\
Als projectnaam is gekozen voor ``CalZone'', ge\"{i}spireerd op het feit dat we een zone voor een kalender moeten maken. Het logo is afgebeeld in figuur \ref{fig:logoProject}.
\begin{figure} [H]
    \centering
    \includegraphics[width = 0.75\textwidth]{Overview/logo_Green_Crop.jpg}
    \caption{Het logo.}
    \label{fig:logoProject}
\end{figure}
%\subsection{Onderstellingen en beperkingen}

\subsection{Project deliverables}
De deliverables voor dit project zijn de volgende:
\begin{itemize}
	\item Software Project Management Plan (SPMP)
	\item Software Test Plan (STD)
	\item Software Requirements Specification (SRS)
	\item Software Design Document (SDD)
	\item Minutes van alle vergaderingen
\end{itemize}
Deze documenten worden in .pdf formaat samen in een enkele zipfile verstuurd via mail. De documenten zijn ook beschikbaar op de GitHub repository. Op het einde van elke iteratie wordt ook de code van het project opgeleverd, deze zit dan in dezelfde zipfile als de documenten. In tabel \ref{tab:kalender} staat de lijst van data waarop de code en de documenten worden opgeleverd. Na elke iteratie volgt een presentatie waar het geleverde werk voor die iteratie wordt gedemonstreert en besproken. 
\begin{table}[H]
  \centering
  \caption{Kalender}
    \begin{tabular}{c|c}
    \textbf{Datum} & \textbf{To Do} \\
    \hline
    Maandag 04/11/2013 & Inleveren SPMP \\
    Vrijdag 15/11/2013 & Eerste versie documenten \\
    Vrijdag 13/12/2013 & Einde iteratie 1: opleveren code en documenten \\
    Woensdag 18/12/2013 & Presentatie \\
    \hline
    \hline
    Dinsdag 04/03/2014 & Einde iteratie 2: opleveren code en documenten \\
    Woensdag 12/03/2014 & Presentatie \\
    Dinsdag 15/04/2014 & Einde iteratie 3: opleveren code en documenten \\
    Woensdag 23/04/2014 & Presentatie \\
    Vrijdag 16/05/2014 & Einde iteratie 4: opleveren code en documenten \\
    Woensdag 21/05/2014 & Finale presentatie \\
    \end{tabular}
  \label{tab:kalender}
\end{table}

Meer gedetailleerde informatie is te vinden in paragraaf \ref{sec:rapportering}.

\section{Evolutie van het SPMP}
De evoluatie van de SPMP zal bijgehouden worden met behulp van een versie geschiedenis in het begin van dit document. Er zullen steeds geplande updates uitgevoerd worden op de tijdsstippen beschreven in tabel \ref{tab:kalender}. 
\\
\\
Deze updates worden uitgevoerd met behulp van de GitHub Repository \cite{GitHubRepository}. Met behulp van GitHub kunnen we met issues werken, en hiervoor telkens een verantwoordelijke aanduiden.